% =============================================================================
% Chapter 3 --- System Architecture
% =============================================================================

\chapter{System Architecture}
\label{ch:architecture}

This chapter presents the controller stack at the block-diagram level:
the three RTL modules, the transaction-word format they agree on, and
the host command set. Module-internal behavior --- FSM state graphs,
pulse-train arithmetic, the LUT-driven programming path --- is deferred
to \Cref{ch:rtl}.

\section{Top-Level Organization}
\label{sec:arch-topology}

The stack sits between a host CPU and the analog Y-Flash core. The host
presents a 32-bit data payload and a multi-bit command token; the
controller decodes that traffic, maintains the closed-loop programming
and verification policy, stages inference pixels, and drives the analog
core over a structured output word and a pulse-mode bus that selects the
active operation (read, program, erase, or inference). The three-module
organization is shown in \Cref{fig:top-level-block-diagram}.

\begin{figure}[h]
    \centering
    \includegraphics[width=\textwidth]{example-image.png}
    \caption{Top-level block diagram. The host drives a transaction
        word and a command token into the combinational
        \emph{Parallel Interface}, which qualifies the packet and
        produces a validated bus consumed by the \emph{ANN Controller}.
        The controller sequences commands toward the analog ANN core
        over a packed core word and a pulse-mode bus, and it
        simultaneously drives the \emph{Input Buffer} that stores
        inference pixels and the expected-weight nibble used during
        verify. A single status bit from the core (\emph{op\_done})
        closes the loop in the program and erase sub-FSMs.}
    \label{fig:top-level-block-diagram}
\end{figure}

\paragraph{Parallel Interface.}
A purely combinational host decoder. It slices the host payload into a
data byte and a one-hot address tail, decodes the tail into a compact
packed address, forwards the command token, and computes a single
\emph{valid} qualifier that asserts only when the command is not idle
\emph{and} every field of the one-hot tail is structurally legal.
Invalid packets never propagate past the interface, so the controller
can assume legal addressing whenever \emph{valid} is high. Lifting all
packet-level validation into a single combinational stage concentrates
the protocol contract in one location and keeps the downstream FSM
simpler.

\paragraph{ANN Controller.}
The only sequential module in the stack. The controller hosts a
nine-state main FSM (idle, reset, program, verify, read, erase, collect,
compute, result) plus three cooperating sub-FSMs that handle programming,
verification, and erase. Two registered captures latch the
host-supplied address and payload on the transition out of idle, so
multi-cycle operations are immune to subsequent host activity. One large
combinational block drives every module output; two clocked blocks
register the FSM states, counters, and captured registers.

\paragraph{Input Buffer.}
A synchronous register-array store that plays two roles. During
inference, it accumulates a pixel window written during \emph{collect}
cycles and exposes the pixels through eight parallel bit-serial lanes
during \emph{compute}. The intended inference workload is MNIST digit
classification, whose $28\times28$ grayscale images are streamed eight
pixels per beat through the buffer, so the window width and lane count
are sized for one MNIST row at a time. During programming, the same
buffer holds the expected-weight nibble for each cell so the verify
comparator can fetch it without an extra host transaction. A single
physical store serving both roles saves area and simplifies addressing
inside the FSM.

\section{Host-to-Core Transaction Word}
\label{sec:arch-word}

The 32-bit transaction word combines a payload byte with four
concatenated one-hot address fields, summarized in
\Cref{tab:ann-core-word-fields} and illustrated in
\Cref{fig:ann-core-word-layout}. The same layout is used at the host
input and the core output; internally the controller operates on a
compact 16-bit packed binary representation of the same address.

\begin{table}[h]
    \centering
    \caption{Fields of the host-to-core transaction word.}
    \label{tab:ann-core-word-fields}
    \begin{tabularx}{\linewidth}{l l l X}
        \toprule
        \textbf{Bits} & \textbf{Width} & \textbf{Field} & \textbf{Meaning} \\
        \midrule
        High byte & 8 & Payload &
            Quantized weight nibble (lower half) on programming, or one
            pixel byte on inference. \\
        Upper nibble  & 4 & \textit{PE} (one-hot) &
            Processing-element / block select. \\
        Lower nibble  & 4 & \textit{SA} (one-hot) &
            Sub-array / sub-block select. \\
        Upper byte (low) & 8 & \textit{col\_ID} (one-hot) &
            Column select within the sub-block. \\
        Lower byte & 8 & \textit{row\_ID} (one-hot) &
            Row select within the sub-block. \\
        \bottomrule
    \end{tabularx}
\end{table}

\begin{figure}[h]
    \centering
    \includegraphics[width=\textwidth]{example-image.png}
    \caption{Pictorial layout of the 32-bit transaction word and the
        internal round-trip used by the controller: the host one-hot
        tail is decoded into a 16-bit packed binary address, the FSM
        captures and operates on the packed address, and the outbound
        core word is rebuilt by single-bit left-shifts into the one-hot
        domain just before reaching the core.}
    \label{fig:ann-core-word-layout}
\end{figure}

\paragraph{Why one-hot at the boundary.}
The one-hot encoding at the host interface mirrors how the analog core
physically selects a cell: only the row and column whose deselect signal
is released conduct, while every other device in the array is held in
its undisturbed state (\Cref{sec:bg-yflash}). Aligning the host-visible
transaction format with the device-level selection mechanism allows the
address to pass through the stack with minimal transformation while
still being structurally validated at the boundary.

\section{Command Set}
\label{sec:arch-commands}

The host command tokens are enumerated in \Cref{tab:cmd-set}. Each
command class drives a specific main-FSM branch and a distinct pulse-mode
encoding for the duration of the active phase.

\begin{table}[h]
    \centering
    \caption{Host command set. Each command class is encoded as a 3-bit
        token on the host \texttt{cmd} input; the same 3-bit value is
        re-used at the core-facing pulse-mode bus, so command decode at
        the boundary and pulse-mode generation at the core share a
        single encoding space.}
    \label{tab:cmd-set}
    \begin{tabularx}{\linewidth}{l c X}
        \toprule
        \textbf{Command} & \textbf{Encoding} & \textbf{Meaning} \\
        \midrule
        \texttt{HIZ}  & \texttt{3'b000} &
            Idle. Host bus is inactive; no transaction starts. The same
            encoding drives the pulse-mode bus to high impedance between
            bursts. \\
        \texttt{READ} & \texttt{3'b001} &
            Read the addressed cell over a read pulse train. \\
        \texttt{PROG} & \texttt{3'b010} &
            Program the addressed cell with the quantized weight in the
            payload; enters the closed-loop program--verify--erase
            policy. \\
        \texttt{ERASE}& \texttt{3'b011} &
            Host-initiated erase of the addressed cell. Distinct from
            the verify-initiated erase taken on an over-programmed
            readback. \\
        \texttt{INF}  & \texttt{3'b100} &
            Inference. Pixel bytes are streamed over consecutive beats;
            the controller stages them in the buffer and then runs a
            compute phase. \\
        \bottomrule
    \end{tabularx}
\end{table}

\section{Module Partitioning Philosophy}
\label{sec:arch-partitioning}

The three-module decomposition is a deliberate architectural partition,
not a code-organization choice. Each module has a single responsibility
with a narrow, structurally validated interface: the interface is
purely combinational and is the only module that touches host-side
encoding; the controller is the only module that contains state, so
every sequential decision in the stack is concentrated in one place;
the input buffer is a synchronous data store with no knowledge of the
high-level operation, exposing a small mode-qualified port set.

The benefit is twofold. First, each module can be verified in isolation
against a focused contract --- one-hot validity for the interface, FSM
transition correctness for the controller, output-and-retention checks
for the buffer --- before integration. Second, replacing or extending
any one module in a future revision (adding a serial host interface,
pipelining the controller, swapping the register array for SRAM) does
not require touching the other two. Bit-level conventions are
centralized in per-module SystemVerilog packages that the testbenches
also reuse, so the protocol is captured in exactly one place.
